\chapter{Conclusion}
\label{chap:conclusion}

This dissertation asked whether crash-triggered Voice traffic can obtain better wireless service than ordinary Best Effort under contention in decentralized vehicular Wi-Fi, and at what cost that protection is paid.
The answer was pursued with a deliberately minimal two-class model: all vehicles exchange periodic Best Effort multicast, one designated crash source injects temporary Voice alert bursts marked DSCP~46, and five \gls{mac} policies are compared on the same highway matrix---plain \gls{dcf}, EDCA with explicit Voice mapping, and three adaptive extensions (\texttt{stable}, \texttt{guarded}, \texttt{emergency}) that suppress competing Best Effort while alerts are active~(\autoref{chap:sysmodel}, \autoref{chap:implementation}).
The evaluation is simulation-only and mechanism-oriented: it isolates statistical contention shaping under controlled load and density, not deterministic real-time guarantees or application-level crash outcomes~(\autoref{chap:evaluation})~\cite{kosekszott2012What,IEEE80211e_2005,aguiar2026_veins_qos}.

%%%%%%%%%%%%%%%%%%%%%%%%%%%%%%%%%%%%%%%%%%%%%%%%%%%%%%%%%%%%%%%%%%%%%%%%%%%%%%%%
\section{Main Findings}
\label{sec:conc:findings}

Under heavy density and high offered load, Voice 95th-percentile alert age is about 2.2~ms under every policy: EDCA does not beat plain \gls{dcf}, and emergency's 2.21~ms versus 2.24~ms is not a material Voice-tail win.
The policy contrast is Best Effort cost.
\texttt{stable} raises fleet Best Effort P95 from 0.51~ms under plain \gls{dcf} to 30~ms; \texttt{guarded} to 3.9~ms; emergency keeps 0.37~ms by dropping competing frames (about $2.5\times10^{3}$ Best Effort drops while blocked)~\cite{aguiar2026_veins_qos}.

The light-density high-load sweep is the same gating test rather than a Voice-gain story.
Bounded adaptive profiles (\texttt{stable}/\texttt{guarded}) leave Voice P95 at 2.08~ms while inflating Best Effort P95 to 49~ms and 9.0~ms; only emergency shortens the Best Effort tail, by preemption rather than by protecting Voice.
Cross-reading both densities shows an enabling condition that this matrix does not meet: adaptive blocking would help Voice only when removing Best Effort competition actually shortens alert service; under a Poisson 125~ms offer it mainly shifts delay onto ordinary traffic.

%%%%%%%%%%%%%%%%%%%%%%%%%%%%%%%%%%%%%%%%%%%%%%%%%%%%%%%%%%%%%%%%%%%%%%%%%%%%%%%%
\section{Practical Takeaway}
\label{sec:conc:takeaway}

The practical takeaway is therefore conditional.
Explicit DSCP-to-Voice mapping remains a reasonable default for crash-marked alerts~\cite{IEEE80211e_2005}.
Adaptive Best Effort suppression---especially emergency preemption---should be treated as an event-triggered escalation level attached to genuine congestion and genuine alert activity, not as an always-on replacement for ordinary EDCA.
Deployment would additionally require authentication, rate limiting, and plausibility checks before overheard Voice frames may arm suppression, because the same local trigger that protects a crash can be abused by faulty or malicious Voice traffic (\autoref{chap:evaluation}).

The absolute Voice delays in this synthetic setting remain far below representative platooning latency budgets such as the~\SI{25}{\milli\second} end-to-end bound in \gls{threegpp} TS~22.186~\cite{ts22186}.
The contribution is accordingly \gls{mac}-level cost accounting under a Poisson offer, not a claim that the simulated scenario already improves crash outcomes at the application layer.

%%%%%%%%%%%%%%%%%%%%%%%%%%%%%%%%%%%%%%%%%%%%%%%%%%%%%%%%%%%%%%%%%%%%%%%%%%%%%%%%
\section{Limitations}
\label{sec:conc:limits}

The study's limits are those already argued in the Evaluation discussion (\autoref{sec:eval:discussion}); they are restated here as the boundary conditions on the Conclusion.

\noindent\textbf{Simulation scope and statistical strength.}
The evaluation is simulation-only, with no testbed or hardware validation, so claims remain tied to the configured Veins/\gls{inet} model~\cite{aguiar2026_veins_qos}.
Heavy-density \glspl{kpi} average three seeds and report the min--max range on headline delay tables, still without confidence intervals or significance tests, and light-density results use the same three-seed rule; the numbers are therefore qualitative mechanism checks under a controlled matrix, not proven effect sizes.
All observed prioritization is statistical contention shaping, not deterministic real-time communication~\cite{kosekszott2012What}.

\medskip
\noindent\textbf{Traffic abstraction and baselines.}
The traffic model is a synthetic two-class Best Effort / Voice abstraction rather than \gls{cam}/\gls{denm} schedules or richer multi-class ITS stacks, and the high crash-burst profile is a deliberate stress case rather than a standardized alert timetable.
Treating \emph{all} Best Effort as suppressible is a worst-case bound for that abstraction: in a real stack, traffic outside the current crash alert may still be safety relevant, so a binary split can underestimate preemption harm (\autoref{sec:eval:disc:cost}).
No adaptive-EDCA or \gls{dcc} baseline is implemented, so no superiority claim is made against those always-on families~\cite{etsi_dcc_2018,romdhani2003AdaptiveEDCF}.
The three adaptive profiles are fixed design points whose triggers are not individually ablated.

\medskip
\noindent\textbf{Radio, topology, and delivery model.}
Connectivity uses free-space path loss, a static \gls{snir} threshold, and binary obstacle loss; small-scale fading, shadowing, and graded attenuation are omitted (\autoref{chap:implementation}).
The corridor is a one-hop highway multicast setting with a single crash source: overlapping protection windows from concurrent alerts are not evaluated, alerts are not retransmitted as multi-hop relays across the 5~km road, and only the light/heavy density pair with three offered-load levels is swept---not a continuous highway-parameter space of speeds, lane layouts, or channel plans.

\medskip
\noindent\textbf{Interpretation caveats that affect how results are read.}
Absolute Voice delays remain far below representative platooning budgets even under plain \gls{dcf} (\autoref{sec:eval:disc:significance})~\cite{ts22186}, so the reported 2.2~ms Voice P95 quantifies alert age under a moderate Poisson offer rather than demonstrated application-level safety benefit, while Best Effort penalties from untimely blocking can still be severe.
Under \texttt{plain}, Best Effort and Voice statistics differ because the traffic processes differ, not because the single \gls{dcf} queue differentiates them (\autoref{sec:eval:disc:plain}); only within-class, cross-policy comparisons are controlled.
Pooled Best Effort averages understate spatially concentrated cost near the crash source (\autoref{sec:eval:disc:spatial}), and unauthenticated Voice frames can arm suppression---a denial-of-service risk bounded by finite windows but not removed without authentication and rate limiting (\autoref{sec:eval:disc:dos}).

Within that scope the comparison remains internally fair---shared radio stack, mobility, and load matrix---so the reported differences isolate \gls{mac}-policy effects~\cite{aguiar2026_veins_qos}.

%%%%%%%%%%%%%%%%%%%%%%%%%%%%%%%%%%%%%%%%%%%%%%%%%%%%%%%%%%%%%%%%%%%%%%%%%%%%%%%%
\section{Future Work}
\label{sec:conc:future}

Several extensions follow directly from the limitations above and from the Evaluation discussion.

\medskip
\noindent\textbf{Stronger evidence and richer scenarios.}
Larger seed batches with confidence intervals or significance tests would move the study from qualitative mechanism checks toward statistically supported effect sizes.
Richer propagation (fading, shadowing, graded obstacles), continuous sweeps of density and offered load, multi-channel or multi-hop relay settings on long corridors, and standardized \gls{cam}/\gls{denm}-like schedules would test whether the conditional frontier survives more realistic radio and traffic stress.
Connecting Voice delay and loss to an application-level safety deadline, crash-avoidance distance, or other outcome metric would address the practical-significance gap left by the present 2.2~ms Voice tails~\cite{ts22186}.

\medskip
\noindent\textbf{Baselines, parameters, and concurrency.}
Activation-criteria ablations and systematic parameter sweeps would separate which trigger and window settings drive the observed Voice gains.
Adaptive-\gls{edca}, ITS-G5 \gls{dcc}, and sidelink baselines would place the event-triggered profiles against established always-on congestion and prioritization families~\cite{etsi_dcc_2018,romdhani2003AdaptiveEDCF}.
Concurrent alert sources---and the interaction of overlapping protection windows---remain unevaluated.

\medskip
\noindent\textbf{Deployment constraints.}
Security against malicious or faulty Voice triggering (authentication, rate limiting, plausibility checks), finer-grained classes than a binary Voice/Best Effort split, and per-vehicle Best Effort fairness around the crash source are prerequisites before comparable \gls{mac} aggression could be defended outside a synthetic two-class study.
Hardware-in-the-loop or testbed validation would further bound how much of the simulated protection-versus-cost frontier transfers to real 802.11p stacks.

In summary, crash-aware prioritization in decentralized vehicular Wi-Fi can measurably improve Voice delay, jitter, and Voice-frame loss under heavy congestion, but only as a carefully bounded, event-triggered escalation whose cost to ordinary traffic must be acknowledged and constrained.
